\documentclass[11pt]{article}
\usepackage[a4paper,margin=27mm]{geometry}
\usepackage{amsmath,amssymb,mathtools}
\usepackage{fontspec,unicode-math}
\setmainfont{Latin Modern Roman}
\setmathfont{Latin Modern Math}
\usepackage{xurl}
\usepackage[hidelinks]{hyperref}
\newcommand{\sref}[2]{\hyperlink{source:#1:#2}{[S#2]}}
\newcommand{\cataloguescope}[1]{\paragraph{Mathematical question.}#1}
\hypersetup{pdftitle={Problem statement}}
\begin{document}
\section{Infinitude of rational points on an elliptic curve}\label{statement}

\subsection{Definitions and mathematical statement}

\paragraph{Shared notation.}$\mathbb N=\{1,2,\ldots\}$, and $\mathbb Z,\mathbb Q,\mathbb R,\mathbb C$ denote the integers, rationals, reals and complex numbers. Any other conventions are specified locally.


Let $A,B\in\mathbb Q$ satisfy $4A^3+27B^2\ne0$, and let $E/\mathbb Q$ be the nonsingular curve $y^2=x^3+Ax+B$ together with its point at infinity $O$. Its rational points form a finitely generated abelian group
\[
E(\mathbb Q)\cong E(\mathbb Q)_{\mathrm{tors}}\oplus\mathbb Z^r,
\]
where the torsion subgroup is finite and $r\in\mathbb Z_{\ge0}$ is the Mordell--Weil rank. Consequently $E(\mathbb Q)$ is infinite exactly when $r>0$.
\paragraph{Statement recorded in the supplied source.}
Given $A,B$ as above, determine whether $E(\mathbb Q)$ is infinite or finite; equivalently, determine whether its Mordell--Weil rank is positive or zero. \sref{statement}{2}
\subsection{Short English statement}
Determine, for an arbitrary rational elliptic curve, whether it has infinitely many rational points.
\subsection{Sources}
\begin{description}
\item[\hypertarget{source:statement:1}{S1}] ULAM.AI and the UnsolvedMath Contributors, UnsolvedMath: A Curated Collection of Open Mathematics Problems, catalogue entry. CC BY 4.0. \url{https://huggingface.co/datasets/ulamai/UnsolvedMath}.
\item[\hypertarget{source:statement:2}{S2}] Christian Wuthrich, \emph{Open problems in number theory}, lecture slides (December 2011) slides headed An elliptic curve and Question, with examples $y^2=x^3+x+1$ and $y^2=x^3+x+2$. \url{https://www.maths.nottingham.ac.uk/plp/pmzcw/download/open.pdf}.
\end{description}
\label{end:statement}
\end{document}
